\documentclass[11pt]{article}
\usepackage[utf8]{inputenc}
\usepackage[T1]{fontenc}
\usepackage{amsmath,amssymb,amsthm}
\usepackage[margin=2.5cm]{geometry}
\usepackage{enumitem}

\newtheorem{theorem}{Theorem}
\newtheorem{lemma}{Lemma}
\theoremstyle{remark}
\newtheorem*{remark}{Remark}

\newcommand{\R}{\mathbb{R}}

\title{The inverse Sturm--Liouville problem:\\
statement, uniqueness, and a reconstruction method}
\author{\texttt{MathPhysicsTheory}}
\date{}

\begin{document}
\maketitle

\begin{abstract}
We summarise what is stated (and, in Lean, proved) about the inverse
Sturm--Liouville problem in this repository: existence and uniqueness up to a
positive constant factor of the coefficients $(p,q,w)$ from the eigenpairs.
We then describe a constructive reconstruction method which turns the
differential equation into a pointwise \emph{linear algebraic} system for
$(p'/p,\,q/p,\,w/p)$, and work out two examples. The reconstruction method is a
mathematical derivation and is \textbf{not} formalised in Lean yet.
\end{abstract}

\section{The problem}

\paragraph{Direct problem.}
Given $p,q,w$ on $[a,b]$ and Robin boundary conditions $B_aX=0$, $B_bX=0$,
find the eigenvalues $\lambda_n$ and eigenfunctions $X_n\neq0$ with
\[
  (p\,X_n')' + q\,X_n + \lambda_n\,w\,X_n = 0 .
\]
The standing assumptions are: $p\in C^1[a,b]$, $p>0$; $q\in C^0[a,b]$,
$q\le 0$; $w\in C^0[a,b]$, $w>0$. Then the eigenvalues satisfy
$0\le\lambda_0<\lambda_1<\cdots$, $\lambda_n\to+\infty$, and $X_n$ has
exactly $n$ zeros in $(a,b)$.

\paragraph{Inverse problem.}
The roles are swapped. Given the interval $[a,b]$, the boundary conditions and
the eigenpairs $(\lambda_n,X_n)_{n\ge0}$, find $(p,q,w)$ satisfying the same
assumptions with
\begin{equation}\label{eq:main}
  (p\,X_n')' + q\,X_n + \lambda_n\,w\,X_n = 0 \quad\text{on }[a,b],\qquad n=0,1,2,\dots
\end{equation}
The data are required to come from \emph{some} Sturm--Liouville problem on the
same interval with the same boundary operators (this is
\texttt{InverseProblem.Assumptions} in the Lean development); in particular
$X_n\in C^2[a,b]$ and $X_n\not\equiv0$.

\section{What is stated and proved}

\begin{theorem}[Existence]
An inverse problem has a solution: the coefficients of any Sturm--Liouville
problem whose eigenpairs are the given data.
\end{theorem}

\begin{theorem}[Uniqueness up to scale]\label{thm:unique}
If $(p_1,q_1,w_1)$ and $(p_2,q_2,w_2)$ both solve the same inverse problem,
then there is a constant $c>0$ with
\[
  p_2=c\,p_1,\qquad q_2=c\,q_1,\qquad w_2=c\,w_1 \quad\text{on }[a,b].
\]
\end{theorem}

The factor $c$ cannot be removed: \eqref{eq:main} is linear in $(p,q,w)$, so
$(cp,cq,cw)$ is a solution whenever $(p,q,w)$ is. Fixing the scale, e.g.\
$p(a)=1$, makes the solution unique. In Lean these are
\texttt{InverseProblem.exists\_solution} and
\texttt{InverseProblem.solution\_unique}; the proof of the latter rests on
\texttt{proportional\_of\_eigenpairs} in
\texttt{InverseProblem/Uniqueness.lean}, which already needs only
\emph{three} eigenpairs. The key tools are: elimination of $y''$ between two
equations, the Riccati equation for the logarithmic flux $m=py'/y$, and a
unique-continuation argument (a solution with $y(x_1)=y'(x_1)=0$ vanishes
identically).

\section{A reconstruction method}

The idea is that \eqref{eq:main} looks like a differential equation for $p$ but
is in fact a linear algebraic equation for the values of the coefficients at a
point.

\paragraph{Step 1: expand.}
Since $p\in C^1$ and $X_n\in C^2$, $(pX_n')'=pX_n''+p'X_n'$, so
$pX_n''+p'X_n'+qX_n+\lambda_nwX_n=0$.

\paragraph{Step 2: divide by $p>0$ and rename.}
Put
\[
  \alpha=\frac{p'}{p},\qquad \beta=\frac{q}{p},\qquad \gamma=\frac{w}{p}.
\]
Then for all $n$ and all $x$
\begin{equation}\label{eq:lin}
  X_n''+\alpha\,X_n'+\beta\,X_n+\lambda_n\gamma\,X_n=0 .
\end{equation}
For fixed $x$ the numbers $X_n(x),X_n'(x),X_n''(x),\lambda_n$ are known data,
so \eqref{eq:lin} is \emph{linear} in the three unknowns
$\alpha(x),\beta(x),\gamma(x)$. Nothing has to be integrated yet.

\paragraph{Step 3: three eigenpairs give a linear system.}
Using $n=0,1,2$, at every $x$
\begin{equation}\label{eq:system}
  \underbrace{\begin{pmatrix}
    X_0' & X_0 & \lambda_0X_0\\
    X_1' & X_1 & \lambda_1X_1\\
    X_2' & X_2 & \lambda_2X_2
  \end{pmatrix}}_{M(x)}
  \begin{pmatrix}\alpha\\ \beta\\ \gamma\end{pmatrix}
  =
  \underbrace{\begin{pmatrix}-X_0''\\ -X_1''\\ -X_2''\end{pmatrix}}_{r(x)} .
\end{equation}

\paragraph{Step 4: Cramer's rule.}
Let $D=\det M$ and let $A,B,G$ be $D$ with the first, second, third column
replaced by $r$. Then
\[
  \alpha=\frac AD,\qquad \beta=\frac BD,\qquad \gamma=\frac GD
  \qquad\text{wherever }D(x)\ne0 .
\]
These hold for \emph{any} solution of the inverse problem, so they are
necessary conditions: the solution cannot take other values.

\paragraph{Step 5: $D\ne0$ on a dense set.}
$D$ may vanish at isolated points (always at a boundary point with Dirichlet
conditions, where all $X_n=0$) but never on an interval. Where $X_n\ne0$,
\[
  D=X_0X_1X_2\det\begin{pmatrix}u_0&1&\lambda_0\\u_1&1&\lambda_1\\u_2&1&\lambda_2\end{pmatrix},
  \qquad u_n=\frac{X_n'}{X_n}.
\]
So $D=0$ on an interval means that the points $(\lambda_n,u_n)$ are collinear.
The functions $m_n=p\,u_n$ satisfy the Riccati equation
$m_n'=-(q+\lambda_nw)-m_n^2/p$. Differentiating the collinearity relation of
the $m_n$ and using this equation gives $m_1=m_0$ on the interval and then
$(\lambda_1-\lambda_0)w=0$, contradicting $w>0$. (This is the argument behind
Theorem~\ref{thm:unique}.) Hence $\{D\ne0\}$ is dense in $(a,b)$, and the
continuous functions $\alpha,\beta,\gamma$ are determined by their values
there; at a zero of $D$ take the limit of $A/D$, $B/D$, $G/D$.

\paragraph{Step 6: recover $p$.}
Since $\alpha=(\ln p)'$,
\[
  p(x)=p(a)\exp\Big(\int_a^x\alpha(t)\,dt\Big).
\]
The constant $p(a)>0$ is exactly the scale freedom of Theorem~\ref{thm:unique};
normalise, e.g., $p(a)=1$.

\paragraph{Step 7: recover $q$ and $w$.}
$q=p\,\beta$ and $w=p\,\gamma$.

\paragraph{Step 8: check.}
Every further eigenpair $(\lambda_n,X_n)$, $n\ge3$, must satisfy
\eqref{eq:lin} with the computed $\alpha,\beta,\gamma$ (automatic for genuine
data; a useful test for noisy data). Also verify $p>0$, $w>0$, $q\le0$.

\begin{remark}[Shortcut]
By Cramer's rule, $(\alpha,\beta,\gamma)$ are the coordinates of $r$ in the
basis formed by the three columns of $M$. If $-X_n''=c_1X_n'+c_2X_n+c_3\lambda_nX_n$
with the \emph{same} $c_1,c_2,c_3$ for $n=0,1,2$, then
$(\alpha,\beta,\gamma)=(c_1,c_2,c_3)$, and one only needs $D\ne0$.
\end{remark}

\begin{remark}[Numerics]
The formulas are ill-conditioned near zeros of $D$ (in particular near the
endpoints for Dirichlet conditions). Normalise $p$ at an interior point, or
fit $(\alpha,\beta,\gamma)$ by least squares from more than three eigenpairs.
A numerical run on $p=1+x^2$, $q=-x$, $w=2+\sin3x$ on $[0,1]$ recovered
$\alpha,\beta,\gamma$ to about $10^{-9}$ and $p,q,w$ to about
$10^{-6}$--$10^{-8}$ away from the endpoints.
\end{remark}

\section{Example 1: the string, $p=1$, $q=0$, $w=1$}

\paragraph{Data.} $X''+\lambda X=0$, $X(0)=X(\pi)=0$, with
$\lambda_n=(n+1)^2$, $X_n=\sin((n+1)x)$. We pretend not to know the equation
and use $n=0,1,2$.

\paragraph{System.} With $X_n'=k\cos kx$, $X_n''=-k^2\sin kx$, $k=n+1$,
\[
  \begin{pmatrix}
    \cos x & \sin x & \sin x\\
    2\cos2x & \sin2x & 4\sin2x\\
    3\cos3x & \sin3x & 9\sin3x
  \end{pmatrix}
  \begin{pmatrix}\alpha\\ \beta\\ \gamma\end{pmatrix}
  =
  \begin{pmatrix}\sin x\\ 4\sin2x\\ 9\sin3x\end{pmatrix}.
\]

\paragraph{Solution.} The right-hand side equals the third column of $M$, so
$A$ and $B$ each have two equal columns and $G=D$:
\[
  A=0,\quad B=0,\quad G=D,\qquad D(x)=-16\sin^6x .
\]
$D\neq0$ on $(0,\pi)$ and vanishes at the endpoints, as predicted in Step~5.
Hence $\alpha=\beta=0$, $\gamma=1$ on $(0,\pi)$, and by continuity on
$[0,\pi]$. Then $p=p(0)\exp\int_0^x0=p(0)$; with $p(0)=1$,
\[
  \boxed{p\equiv1,\quad q\equiv0,\quad w\equiv1.}
\]
\paragraph{Check.} For $n=3$: $X_3''+16X_3=-16\sin4x+16\sin4x=0$.

\section{Example 2: variable coefficients, $p=x$, $q=0$, $w=1/x$}

\paragraph{Data.} On $[1,e^\pi]$, $(xX')'+\frac{\lambda}{x}X=0$,
$X(1)=X(e^\pi)=0$. With $t=\ln x$ this becomes $X_{tt}+\lambda X=0$ on
$[0,\pi]$, so $\lambda_n=(n+1)^2$, $X_n=\sin((n+1)\ln x)$.

\paragraph{Derivatives.} With $k=n+1$,
\[
  X_n'=\frac{k\cos(k\ln x)}{x},\qquad
  X_n''=-\frac{k^2\sin(k\ln x)}{x^2}-\frac{k\cos(k\ln x)}{x^2}
       =-\frac{\lambda_nX_n}{x^2}-\frac{X_n'}{x}.
\]

\paragraph{Solution.} For $n=0,1,2$
\[
  -X_n''=\frac1x\,X_n'+0\cdot X_n+\frac1{x^2}\,\lambda_nX_n
\]
with the same coefficients in all three rows, so by the shortcut
\[
  \alpha=\frac1x,\qquad\beta=0,\qquad\gamma=\frac1{x^2}.
\]
(Concretely $A=D/x$, $B=0$, $G=D/x^2$ with $D=-16\sin^6(\ln x)/x$, nonzero on
$(1,e^\pi)$.) Then $p=p(1)\exp\int_1^x\frac{dt}{t}=p(1)\,x$, and with $p(1)=1$,
\[
  \boxed{p(x)=x,\quad q\equiv0,\quad w(x)=\frac1x,}
\]
the original coefficients up to the scale $p(1)$.

\paragraph{Check.} For $n=3$: $X''+\frac1xX'+\frac{16}{x^2}X=0$ by the formula
for $X_n''$ above.

\section{Summary of the algorithm}
\begin{enumerate}[nosep]
  \item Take $(\lambda_n,X_n)$, $n=0,1,2$, and compute $X_n',X_n''$.
  \item Form $M(x)$ and $r(x)$ as in \eqref{eq:system}.
  \item Solve $M(\alpha,\beta,\gamma)^{T}=r$ (Cramer, or read off coordinates).
  \item $p=p(a)\exp\int_a^x\alpha$, $q=p\beta$, $w=p\gamma$.
  \item Check with $n\ge3$ and the signs $p,w>0$, $q\le0$; the scale $p(a)$ is free.
\end{enumerate}

\paragraph{Open items.} Formalising the reconstruction in Lean: (i) the
necessary conditions of Step~4 (essentially available from the elimination
lemmas in \texttt{Uniqueness.lean}), (ii) density of $\{D\neq0\}$ and the
limit extension at zeros of $D$, (iii) the integral formula for $p$, and
(iv) sufficiency: that the reconstructed triple really solves \eqref{eq:main}
for \emph{all} $n$, which uses the data coming from a genuine
Sturm--Liouville problem.

\end{document}
